\documentclass[12pt]{amsart}
\usepackage{ amsmath, amsthm, amsfonts, amssymb, color}
 \usepackage{mathrsfs}
\usepackage{amsfonts, amsmath}
 \usepackage{amsmath,amstext,amsthm,amssymb,amsxtra}
 \usepackage{txfonts} %also pxfonts
 \usepackage[colorlinks, citecolor=blue,pagebackref,hypertexnames=false]{hyperref}
 \allowdisplaybreaks
 \usepackage{pgf,tikz}
 \usepackage{multirow}%newtablepackage1
 \usepackage{diagbox} %newtablepackage2
\usepackage{cite}
 \usepackage{tikz}

\usetikzlibrary{decorations.pathreplacing}

 %\usepackage{showkeys}
 \textwidth =166mm
 \textheight =232mm
\marginparsep=0cm
\oddsidemargin=2mm
\evensidemargin=2mm
\headheight=13pt
\headsep=0.8cm
\parskip=0pt
\hfuzz=6pt
\widowpenalty=10000
 \setlength{\topmargin}{-0.2cm}

\DeclareMathOperator*{\essinf}{ess\ inf}
\DeclareMathOperator*{\esssup}{ess\ sup}
\setcounter {tocdepth}{1}
\begin{document}

 \baselineskip 16.6pt
\hfuzz=6pt

\widowpenalty=10000

\newtheorem{cl}{Claim}
\newtheorem{theorem}{Theorem}[section]
\newtheorem{proposition}[theorem]{Proposition}
\newtheorem{coro}[theorem]{Corollary}
\newtheorem{lemma}[theorem]{Lemma}
\newtheorem{definition}[theorem]{Definition}
\newtheorem{nota}[theorem]{Notation}
%\theoremstyle{definition}
\newtheorem{assum}{Assumption}[section]
\newtheorem{example}[theorem]{Example}
\newtheorem{remark}[theorem]{Remark}
\renewcommand{\theequation}
{\thesection.\arabic{equation}}

\def\SL{\sqrt H}

\newcommand{\mar}[1]{{\marginpar{\sffamily{\scriptsize
        #1}}}}

\newcommand{\as}[1]{{\mar{AS:#1}}}
\newcommand{\cE}{\mathcal{E}}
\newcommand\R{\mathbb{R}}
\newcommand\RR{\mathbb{R}}
\newcommand\CC{\mathbb{C}}
\newcommand\D {{\rm Dunkl}}
\newcommand\NN{\mathbb{N}}
\newcommand\N {{\bf N}}
\newcommand\ZZ{\mathbb{Z}}
\newcommand\HH{\mathbb{H}}
\newcommand\Z{\mathbb{Z}}
\def\RN {\mathbb{R}^n}
\renewcommand\Re{\operatorname{Re}}
\renewcommand\Im{\operatorname{Im}}

\newcommand{\mc}{\mathcal}

\def\hs{\hspace{0.33cm}}
\newcommand{\la}{\alpha}
\def \l {\alpha}
\newcommand{\eps}{\tau}
\newcommand{\pl}{\partial}
\newcommand{\supp}{{\rm supp}{\hspace{.05cm}}}
\newcommand{\x}{\times}
\newcommand{\lag}{\langle}
\newcommand{\rag}{\rangle}

\newcommand\wrt{\,{\rm d}}

\newcommand{\norm}[2]{|#1|_{#2}}
\newcommand{\Norm}[2]{\|#1\|_{#2}}

\title[Schatten class properties of Dunkl--Riesz transform commutators]{Weyl--Chamber geometry, Poincar\'e inequality meets Schatten class properties of Dunkl--Riesz transform commutators}



\author{Zhijie Fan}
\address{Zhijie Fan, School of Mathematical Sciences, South China Normal University, Guangzhou 510631, China}
\email{fanzhj3@mail2.sysu.edu.cn}







\author{Ji Li}
\address{Ji Li, Department of Mathematics, Macquarie University,NSW2109, Australia}
\email{ji.li@mq.edu.au}



\author{Dmitriy Zanin}
\address{Dmitriy Zanin, University of New South Wales, Kensington, NSW 2052, Australia}
\email{d.zanin@unsw.edu.au}






  \date{\today}

 \subjclass[2010]{47B10, 42B20, 43A85}
\keywords{Schatten class, Riesz transform commutators, Dunkl operator, Besov space, Nearly weakly orthogonal sequence}




\begin{abstract}

\end{abstract}

\maketitle


\tableofcontents

%\newpage

\section{Preliminaries}
\setcounter{equation}{0}
\subsection{Harmonic analysis in the Dunkl setting}
Let $\mathbb{R}^N$ be the Euclidean space equipped with the usual scalar product $\langle \cdot,\,\cdot\rangle $ and the induced norm $\|\cdot\|$. For a nonzero vector $v\in \mathbb{R}^N$, the reflection $\sigma_v$ with respect to the hyperplane $v^{\perp}$ orthogonal to $v$ is defined by
\begin{align*}
	\sigma_v x:= x-2\frac{\langle x,v\rangle}{\|v\|^2}v.
\end{align*}
A finite set $R\subset \mathbb{R}^N\backslash \{0\}$ is called a root system if $\sigma_v(R)=R$ for every $v\in R$. A root system $R$ is called a normalized reduced root system if  $\|v\|^2=2$ for every $v\in R$.
The finite group $G$ generated by the reflections $\{\sigma_v\}_{v\in R}$ is called the Weyl group (reflection group) of the root system $R$.
Let $k: R \to [0,\infty)$ be a multiplicity function, which is a $G$-invariant function  fixed throughout the paper. 

Define the Dunkl measure $dw$ on $\mathbb{R}^N$ by
$$
dw(x) :=\prod_{v \in R} |\langle x,v\rangle|^{k(v)} dx,
$$
where $dx$ stands for the Lebesgue measure on $\mathbb{R}^N$. Let $B(x,r):=\{y\in \mathbb{R}^N:\|x-y\|< r\}$ be the Euclidean ball centred at $x$ with radius $r>0$. Let $\N:=N+\sum_{v\in R}k(v)$. Note that
$$w(B(tx,tr))=t^{\N}w(B(x,r)),\ {\rm for}\ x\in\mathbb{R}^N,\ t,r>0.$$
Furthermore,
\begin{align}\label{twosideinequ}
  w(B(x,r))\sim r^N \prod_{v \in R}(|\langle x, v \rangle |+r)^{k(v)},
\end{align}
where the implicit constant only depends on $N,R$ and $k$.  It follows from \eqref{twosideinequ} that there exists a constant $C\geq 1$ such that for every $x\in \mathbb{R}^N$,
\begin{align}\label{doublinginequality}
C^{-1}\left(\frac{r_2}{r_1}\right)^N\leq \frac{w(B(x,r_2))}{w(B(x,r_1))}\leq C\left(\frac{r_2}{r_1}\right)^\N,\ {\rm for}\ 0<r_1<r_2,
\end{align}
so the numbers $\N$ and $N$ are called the upper and lower homogeneous dimensions, respectively.


\begin{definition}{\rm
The Weyl chambers are defined as the closures of connected components of the set
  \begin{align}
    \left\{x\in\mathbb{R}^N:\langle x,v\rangle \neq 0\ \ { for} \ { all} \ v\in R\right\}.\end{align}}
\end{definition}
Define $
	d(x,y):=\min_{\sigma\in G}\|x-\sigma(y)\|
$
as the distance between two $G$-orbits $\mathcal O(x)$ and $\mathcal O(y)$. Recall from \cite[Chapter VII, proof of Theorem 2.12]{MR514561} that
\begin{align}\label{distanceequa}
d(x,y)=\|x-y\|\ {\rm if}\ {\rm and}\ {\rm only}\ {\rm if}\ x,y\in\Gamma
\end{align}
for some closed Weyl chamber $\Gamma$.


\begin{definition}{\rm\label{beso}
Suppose $1< p< \infty$. We say that a  function $f\in L^p_{{\rm loc}}(\mathbb{R}^{N},dw)$ belongs to {\color{red}Dunkl} Besov space $B^{p}_{\D}(\mathbb{R}^N)$ if
\begin{align*}
\|f\|_{B^{p}_{\D}(\mathbb{R}^N)}:=\left(\int_{\mathbb{R}^{N}}\int_{\mathbb{R}^{N}}\left(\frac{d(x,y)}{\|x-y\|}\right)^{p}\frac{|f(x)-f(y)|^p}{w(B(x,d(x,y)))^2}dw(x)dw(y)\right)^{1/p}<+\infty.
\end{align*}}
\end{definition}

\begin{definition}{\rm\label{beso}
Suppose $1< p< \infty$. We say that a  function $f\in L^p_{{\rm loc}}(\mathbb{R}^{N},dw)$ belongs to Besov space $B^p(\mathbb{R}^N,dw)$ if
$$\|f\|_{B^p(\mathbb{R}^N,dw)}:=\Bigg({\int_{\mathbb{R}^N}\int_{\mathbb{R}^N}\frac{|f(x)-f(y)|^p}{w(B(x,\|x-y\|))^2}}dw(x)dw(y)\Bigg)^{1/p}$$}
\end{definition}
\section{Equivalence of Besov spaces}
\begin{lemma}\label{lem1}
We have
\begin{align*}
&\|f\|_{B^{p}_{\D}(\mathbb{R}^N)}\\
&\sim  \sum_{\sigma\in G} \left(\int_{\sigma(\Gamma)}\int_{\sigma(\Gamma)}\frac{|f(x)-f(y)|^p}{w(B(x,\|x-y\|))^2}dw(x)dw(y)\right)^{1/p}\\
&+\sum_{\substack{\sigma,\tilde{\sigma}\in G\\ \sigma\neq \tilde{\sigma}}}\left(\int_{\Gamma}\int_{\Gamma}\left(\frac{\|x-y\|}{\|\sigma(x)-\tilde{\sigma}(y)\|}\right)^{p}\frac{|f(\sigma(x))-f(\tilde{\sigma}(y))|^p}{w(B(\sigma(x),\|x-y\|))^2}dw(x)dw(y)\right)^{1/p}.
\end{align*}
\end{lemma}
\begin{proof}
Fix a chamber $\Gamma$.
\begin{align*}
&\|f\|_{B^{p}_{\D}(\mathbb{R}^N)}\\
&=\left(\int_{\mathbb{R}^{N}}\int_{\mathbb{R}^{N}}\left(\frac{d(x,y)}{\|x-y\|}\right)^{p}\frac{|f(x)-f(y)|^p}{w(B(x,d(x,y)))^2}dw(x)dw(y)\right)^{1/p}\\
&\sim \sum_{\sigma\in G}\sum_{\tilde{\sigma}\in G}\left(\int_{\sigma(\Gamma)}\int_{\tilde{\sigma}(\Gamma)}\left(\frac{d(x,y)}{\|x-y\|}\right)^{p}\frac{|f(x)-f(y)|^p}{w(B(x,d(x,y)))^2}dw(x)dw(y)\right)^{1/p}\\
&\sim  \sum_{\sigma\in G} \left(\int_{\sigma(\Gamma)}\int_{\sigma(\Gamma)}\left(\frac{d(x,y)}{\|x-y\|}\right)^{p}\frac{|f(x)-f(y)|^p}{w(B(x,d(x,y)))^2}dw(x)dw(y)\right)^{1/p}\\
&+\sum_{\substack{\sigma,\tilde{\sigma}\in G\\\sigma\neq \tilde{\sigma}}}\left(\int_{\sigma(\Gamma)}\int_{\tilde{\sigma}(\Gamma)}\left(\frac{d(x,y)}{\|x-y\|}\right)^{p}\frac{|f(x)-f(y)|^p}{w(B(x,d(x,y)))^2}dw(x)dw(y)\right)^{1/p}\\
&\sim  \sum_{\sigma\in G} \left(\int_{\sigma(\Gamma)}\int_{\sigma(\Gamma)}\left(\frac{d(x,y)}{\|x-y\|}\right)^{p}\frac{|f(x)-f(y)|^p}{w(B(x,d(x,y)))^2}dw(x)dw(y)\right)^{1/p}\\
&+\sum_{\substack{\sigma,\tilde{\sigma}\in G\\\sigma\neq \tilde{\sigma}}}\left(\int_{\Gamma}\int_{\Gamma}\left(\frac{d(\sigma(x),\tilde{\sigma}(y))}{\|\sigma(x)-\tilde{\sigma}(y)\|}\right)^{p}\frac{|f(\sigma(x))-f(\tilde{\sigma}(y))|^p}{w(B(\sigma(x),d(\sigma(x),\tilde{\sigma}(y))))^2}dw(x)dw(y)\right)^{1/p}\\
&\sim  \sum_{\sigma\in G} \left(\int_{\sigma(\Gamma)}\int_{\sigma(\Gamma)}\frac{|f(x)-f(y)|^p}{w(B(x,\|x-y\|))^2}dw(x)dw(y)\right)^{1/p}\\
&+\sum_{\substack{\sigma,\tilde{\sigma}\in G\\ \sigma\neq \tilde{\sigma}}}\left(\int_{\Gamma}\int_{\Gamma}\left(\frac{\|x-y\|}{\|\sigma(x)-\tilde{\sigma}(y)\|}\right)^{p}\frac{|f(\sigma(x))-f(\tilde{\sigma}(y))|^p}{w(B(\sigma(x),\|x-y\|))^2}dw(x)dw(y)\right)^{1/p}.
\end{align*}
\end{proof}

\begin{lemma}\label{lem2}
We have
\begin{align*}
&\|f\|_{B^p(\mathbb{R}^N,dw)}\\
&\sim\sum_{\sigma\in G}\Bigg({\int_{\sigma(\Gamma)}\int_{{\sigma}(\Gamma)}\frac{|f(x)-f(y)|^p}{w(B(x,\|x-y\|))^2}}dw(x)dw(y)\Bigg)^{1/p}\\
&+\sum_{\substack{\sigma,\tilde{\sigma}\in G\\ \sigma\neq \tilde{\sigma}}}\Bigg({\int_{\Gamma}\int_{\Gamma}\frac{|f(\sigma(x))-f(\tilde{\sigma}(y))|^p}{w(B(\sigma(x),\|\sigma(x)-\tilde{\sigma}(y)\|))^2}}dw(x)dw(y)\Bigg)^{1/p}
\end{align*}
\end{lemma}
\begin{proof}
\begin{align*}
&\|f\|_{B^p(\mathbb{R}^N,dw)}\\
&=\Bigg({\int_{\mathbb{R}^N}\int_{\mathbb{R}^N}\frac{|f(x)-f(y)|^p}{w(B(x,\|x-y\|))^2}}dw(x)dw(y)\Bigg)^{1/p}\\
&\sim\sum_{\sigma\in G}\sum_{\tilde{\sigma}\in G}\Bigg({\int_{\sigma(\Gamma)}\int_{\tilde{\sigma}(\Gamma)}\frac{|f(x)-f(y)|^p}{w(B(x,\|x-y\|))^2}}dw(x)dw(y)\Bigg)^{1/p}\\
&\sim\sum_{\sigma\in G}\Bigg({\int_{\sigma(\Gamma)}\int_{{\sigma}(\Gamma)}\frac{|f(x)-f(y)|^p}{w(B(x,\|x-y\|))^2}}dw(x)dw(y)\Bigg)^{1/p}\\
&+\sum_{\substack{\sigma,\tilde{\sigma}\in G\\ \sigma\neq \tilde{\sigma}}}\Bigg({\int_{\sigma(\Gamma)}\int_{\tilde{\sigma}(\Gamma)}\frac{|f(x)-f(y)|^p}{w(B(x,\|x-y\|))^2}}dw(x)dw(y)\Bigg)^{1/p}\\
&\sim\sum_{\sigma\in G}\Bigg({\int_{\sigma(\Gamma)}\int_{{\sigma}(\Gamma)}\frac{|f(x)-f(y)|^p}{w(B(x,\|x-y\|))^2}}dw(x)dw(y)\Bigg)^{1/p}\\
&+\sum_{\substack{\sigma,\tilde{\sigma}\in G\\ \sigma\neq \tilde{\sigma}}}\Bigg({\int_{\Gamma}\int_{\Gamma}\frac{|f(\sigma(x))-f(\tilde{\sigma}(y))|^p}{w(B(\sigma(x),\|\sigma(x)-\tilde{\sigma}(y)\|))^2}}dw(x)dw(y)\Bigg)^{1/p}
\end{align*}
\end{proof}
Define the measure $\mu$ on $\Gamma\times \Gamma$ by setting
$$d\mu(x,y)=\frac{1}{w(B(x,\|x-y\|))^2}dw(x)dw(y).$$
Let
$$F_{\sigma}(x,y)=f(\sigma(x))-f(\sigma(y)).$$
$$F_{\sigma,\tilde{\sigma}}^{(1)}(x,y)=\frac{\|x-y\|}{\|\sigma(x)-\tilde{\sigma}(y)\|}\cdot(f(\sigma(x))-f(\tilde{\sigma}(y))).$$
$$F_{\sigma,\tilde{\sigma}}^{(2)}(x,y)=\frac{\|x-y\|}{\|\sigma(x)-\tilde{\sigma}(y)\|}\cdot(f(\sigma(x))-f(\sigma(y))).$$
$$F_{\sigma,\tilde{\sigma}}^{(3)}(x,y)=\frac{\|x-y\|}{\|\sigma(x)-\tilde{\sigma}(y)\|}\cdot(f(\sigma(y))-f(\tilde{\sigma}(y))).$$
\begin{lemma}
We have
$$\|f\|_{B^{p}_{\D}(\mathbb{R}^N)}\sim \sum_{\sigma\in G} \|F_{\sigma}\|_{L^p(\mu)}+\sum_{\substack{\sigma,\tilde{\sigma}\in G\\ \sigma\neq \tilde{\sigma}}}\|F_{\sigma,\tilde{\sigma}}^{(3)}\|_{L^p(\mu)}$$
\end{lemma}
\begin{proof}
Note that
$F_{\sigma,\tilde{\sigma}}^{(1)}(x,y)=F_{\sigma,\tilde{\sigma}}^{(2)}(x,y)+F_{\sigma,\tilde{\sigma}}^{(3)}(x,y).$
By Lemma \ref{lem1},
We have
\begin{align}\label{susti}
&\|f\|_{B^{p}_{\D}(\mathbb{R}^N)}\nonumber\\
&\sim  \sum_{\sigma\in G} \left(\int_{\sigma(\Gamma)}\int_{\sigma(\Gamma)}\frac{|f(x)-f(y)|^p}{w(B(x,\|x-y\|))^2}dw(x)dw(y)\right)^{1/p}\nonumber\\
&+\sum_{\substack{\sigma,\tilde{\sigma}\in G\\ \sigma\neq \tilde{\sigma}}}\left(\int_{\Gamma}\int_{\Gamma}\left(\frac{\|x-y\|}{\|\sigma(x)-\tilde{\sigma}(y)\|}\right)^{p}\frac{|f(\sigma(x))-f(\tilde{\sigma}(y))|^p}{w(B(\sigma(x),\|x-y\|))^2}dw(x)dw(y)\right)^{1/p}\nonumber\\
&\sim  \sum_{\sigma\in G} \left(\int_{\Gamma}\int_{\Gamma}\frac{|f(\sigma(x))-f(\sigma(y))|^p}{w(B(x,\|x-y\|))^2}dw(x)dw(y)\right)^{1/p}\nonumber\\
&+\sum_{\substack{\sigma,\tilde{\sigma}\in G\\ \sigma\neq \tilde{\sigma}}}\left(\int_{\Gamma}\int_{\Gamma}\left(\frac{\|x-y\|}{\|\sigma(x)-\tilde{\sigma}(y)\|}\right)^{p}\frac{|f(\sigma(x))-f(\tilde{\sigma}(y))|^p}{w(B(x,\|x-y\|))^2}dw(x)dw(y)\right)^{1/p}\nonumber\\
&\sim  \sum_{\sigma\in G} \|F_{\sigma}\|_{L^p(\mu)}+\sum_{\substack{\sigma,\tilde{\sigma}\in G\\ \sigma\neq \tilde{\sigma}}}\|F_{\sigma,\tilde{\sigma}}^{(1)}\|_{L^p(\mu)}
\end{align}
We claim that
\begin{align}\label{claim1}
\sum_{\sigma\in G} \|F_{\sigma}\|_{L^p(\mu)}+\sum_{\substack{\sigma,\tilde{\sigma}\in G\\ \sigma\neq \tilde{\sigma}}}\|F_{\sigma,\tilde{\sigma}}^{(1)}\|_{L^p(\mu)}\sim \sum_{\sigma\in G} \|F_{\sigma}\|_{L^p(\mu)}+\sum_{\substack{\sigma,\tilde{\sigma}\in G\\ \sigma\neq \tilde{\sigma}}}\|F_{\sigma,\tilde{\sigma}}^{(3)}\|_{L^p(\mu)}.
\end{align}
Note that
$$\|F_{\sigma,\tilde{\sigma}}^{(1)}\|_{L^p(\mu)}\leq \|F_{\sigma,\tilde{\sigma}}^{(2)}\|_{L^p(\mu)}+\|F_{\sigma,\tilde{\sigma}}^{(3)}\|_{L^p(\mu)}\leq \|F_{\sigma}\|_{L^p(\mu)}+\|F_{\sigma,\tilde{\sigma}}^{(3)}\|_{L^p(\mu)}.$$
This implies that
$$\sum_{\sigma\in G} \|F_{\sigma}\|_{L^p(\mu)}+\sum_{\substack{\sigma,\tilde{\sigma}\in G\\ \sigma\neq \tilde{\sigma}}}\|F_{\sigma,\tilde{\sigma}}^{(1)}\|_{L^p(\mu)}\lesssim \sum_{\sigma\in G} \|F_{\sigma}\|_{L^p(\mu)}+\sum_{\substack{\sigma,\tilde{\sigma}\in G\\ \sigma\neq \tilde{\sigma}}}\|F_{\sigma,\tilde{\sigma}}^{(3)}\|_{L^p(\mu)}.$$
For the reverse inequality, note that
\[
F_{\sigma,\tilde{\sigma}}^{(3)}
=
F_{\sigma,\tilde{\sigma}}^{(1)}-F_{\sigma,\tilde{\sigma}}^{(2)}.
\]
Hence, by the triangle inequality,
\[
\|F_{\sigma,\tilde{\sigma}}^{(3)}\|_{L^p(\mu)}
\le
\|F_{\sigma,\tilde{\sigma}}^{(1)}\|_{L^p(\mu)}
+
\|F_{\sigma,\tilde{\sigma}}^{(2)}\|_{L^p(\mu)}.
\]
Therefore,
\[
\sum_{\sigma\in G} \|F_{\sigma}\|_{L^p(\mu)}
+\sum_{\substack{\sigma,\tilde{\sigma}\in G\\ \sigma\neq \tilde{\sigma}}}
\|F_{\sigma,\tilde{\sigma}}^{(3)}\|_{L^p(\mu)}
\]
\[
\le
\sum_{\sigma\in G} \|F_{\sigma}\|_{L^p(\mu)}
+
\sum_{\substack{\sigma,\tilde{\sigma}\in G\\ \sigma\neq \tilde{\sigma}}}
\|F_{\sigma,\tilde{\sigma}}^{(1)}\|_{L^p(\mu)}
+
\sum_{\substack{\sigma,\tilde{\sigma}\in G\\ \sigma\neq \tilde{\sigma}}}
\|F_{\sigma,\tilde{\sigma}}^{(2)}\|_{L^p(\mu)}.
\]
Using again
\[
\|F_{\sigma,\tilde{\sigma}}^{(2)}\|_{L^p(\mu)}
\le
\|F_\sigma\|_{L^p(\mu)},
\]
we obtain
\[
\sum_{\sigma\in G} \|F_{\sigma}\|_{L^p(\mu)}
+\sum_{\substack{\sigma,\tilde{\sigma}\in G\\ \sigma\neq \tilde{\sigma}}}
\|F_{\sigma,\tilde{\sigma}}^{(3)}\|_{L^p(\mu)}
\le
|G|\sum_{\sigma\in G}\|F_\sigma\|_{L^p(\mu)}
+
\sum_{\substack{\sigma,\tilde{\sigma}\in G\\ \sigma\neq \tilde{\sigma}}}
\|F_{\sigma,\tilde{\sigma}}^{(1)}\|_{L^p(\mu)}.
\]
Consequently,
\[
\sum_{\sigma\in G} \|F_{\sigma}\|_{L^p(\mu)}
+\sum_{\substack{\sigma,\tilde{\sigma}\in G\\ \sigma\neq \tilde{\sigma}}}
\|F_{\sigma,\tilde{\sigma}}^{(3)}\|_{L^p(\mu)}
\le
|G|
\Biggl(
\sum_{\sigma\in G}\|F_\sigma\|_{L^p(\mu)}
+
\sum_{\substack{\sigma,\tilde{\sigma}\in G\\ \sigma\neq \tilde{\sigma}}}
\|F_{\sigma,\tilde{\sigma}}^{(1)}\|_{L^p(\mu)}
\Biggr).
\]
This ends the proof of the claim. Substituting \eqref{claim1} into \eqref{susti}, we deduce that
$$\|f\|_{B^{p}_{\D}(\mathbb{R}^N)}\sim \sum_{\sigma\in G} \|F_{\sigma}\|_{L^p(\mu)}+\sum_{\substack{\sigma,\tilde{\sigma}\in G\\ \sigma\neq \tilde{\sigma}}}\|F_{\sigma,\tilde{\sigma}}^{(3)}\|_{L^p(\mu)}.$$
\end{proof}

Take
$$G_{\sigma,\tilde{\sigma}}^{(1)}(x,y)=f(\sigma(x))-f(\tilde{\sigma}(y)).$$
$$G_{\sigma,\tilde{\sigma}}^{(2)}(x,y)=f(\sigma(x))-f(\sigma(y)).$$
$$G_{\sigma,\tilde{\sigma}}^{(3)}(x,y)=f(\sigma(y))-f(\tilde{\sigma}(y)).$$
Define the measure $\mu_{\sigma,\tilde{\sigma}}$ on $\Gamma\times \Gamma$ by setting
$$d\mu_{\sigma,\tilde{\sigma}}(x,y)=\frac{1}{w(B(x,\|\sigma(x)-\tilde{\sigma}(y)\|))^2}dw(x)dw(y).$$

\begin{lemma}\label{lem:non-dunkl-G3}
We have
\[
\|f\|_{B^p(\mathbb R^N,dw)}
\sim
\sum_{\sigma\in G}\|F_\sigma\|_{L^p(\mu)}
+
\sum_{\substack{\sigma,\tilde{\sigma}\in G\\ \sigma\neq \tilde{\sigma}}}
\|G_{\sigma,\tilde{\sigma}}^{(3)}\|_{L^p(\mu_{\sigma,\tilde{\sigma}})}.
\]
\end{lemma}

\begin{proof}
Note that
\[
G_{\sigma,\tilde{\sigma}}^{(1)}
=
G_{\sigma,\tilde{\sigma}}^{(2)}
+
G_{\sigma,\tilde{\sigma}}^{(3)}.
\]
By Lemma \ref{lem2},
\begin{align}
\|f\|_{B^p(\mathbb R^N,dw)}
&\sim
\sum_{\sigma\in G}
\Bigg(
\int_{\sigma(\Gamma)}\int_{\sigma(\Gamma)}
\frac{|f(x)-f(y)|^p}{w(B(x,\|x-y\|))^2}\,dw(x)\,dw(y)
\Bigg)^{1/p}
\nonumber\\
&\quad+
\sum_{\substack{\sigma,\tilde{\sigma}\in G\\ \sigma\neq \tilde{\sigma}}}
\Bigg(
\int_{\Gamma}\int_{\Gamma}
\frac{|f(\sigma(x))-f(\tilde{\sigma}(y))|^p}
{w(B(\sigma(x),\|\sigma(x)-\tilde{\sigma}(y)\|))^2}
\,dw(x)\,dw(y)
\Bigg)^{1/p}.
\label{eq:non-dunkl-decomp-1}
\end{align}
Since \(dw\) is \(G\)-invariant and \(\sigma(B(x,r))=B(\sigma(x),r)\), we have
\[
w(B(\sigma(x),r))=w(B(x,r))
\]
for every \(\sigma\in G\), \(x\in\mathbb R^N\), and \(r>0\). Hence \eqref{eq:non-dunkl-decomp-1} becomes
\begin{align}
\|f\|_{B^p(\mathbb R^N,dw)}
&\sim
\sum_{\sigma\in G}
\Bigg(
\int_{\Gamma}\int_{\Gamma}
\frac{|f(\sigma(x))-f(\sigma(y))|^p}{w(B(x,\|x-y\|))^2}
\,dw(x)\,dw(y)
\Bigg)^{1/p}
\nonumber\\
&\quad+
\sum_{\substack{\sigma,\tilde{\sigma}\in G\\ \sigma\neq \tilde{\sigma}}}
\Bigg(
\int_{\Gamma}\int_{\Gamma}
\frac{|f(\sigma(x))-f(\tilde{\sigma}(y))|^p}
{w(B(x,\|\sigma(x)-\tilde{\sigma}(y)\|))^2}
\,dw(x)\,dw(y)
\Bigg)^{1/p}
\nonumber\\
&=
\sum_{\sigma\in G}\|F_\sigma\|_{L^p(\mu)}
+
\sum_{\substack{\sigma,\tilde{\sigma}\in G\\ \sigma\neq \tilde{\sigma}}}
\|G_{\sigma,\tilde{\sigma}}^{(1)}\|_{L^p(\mu_{\sigma,\tilde{\sigma}})}.
\label{eq:non-dunkl-decomp-2}
\end{align}

We claim that
\begin{align}
&\sum_{\sigma\in G}\|F_\sigma\|_{L^p(\mu)}
+
\sum_{\substack{\sigma,\tilde{\sigma}\in G\\ \sigma\neq \tilde{\sigma}}}
\|G_{\sigma,\tilde{\sigma}}^{(1)}\|_{L^p(\mu_{\sigma,\tilde{\sigma}})}
\nonumber\\
&\sim
\sum_{\sigma\in G}\|F_\sigma\|_{L^p(\mu)}
+
\sum_{\substack{\sigma,\tilde{\sigma}\in G\\ \sigma\neq \tilde{\sigma}}}
\|G_{\sigma,\tilde{\sigma}}^{(3)}\|_{L^p(\mu_{\sigma,\tilde{\sigma}})}.
\label{eq:claim-non-dunkl}
\end{align}

We first prove the upper bound. By the triangle inequality,
\[
\|G_{\sigma,\tilde{\sigma}}^{(1)}\|_{L^p(\mu_{\sigma,\tilde{\sigma}})}
\le
\|G_{\sigma,\tilde{\sigma}}^{(2)}\|_{L^p(\mu_{\sigma,\tilde{\sigma}})}
+
\|G_{\sigma,\tilde{\sigma}}^{(3)}\|_{L^p(\mu_{\sigma,\tilde{\sigma}})}.
\]
Now for \(x,y\in \Gamma\), since \(x\) and \(y\) belong to the same closed Weyl chamber, we have
\[
d(x,y)=\|x-y\|.
\]
By \(G\)-invariance of the orbit distance,
\[
d(\sigma(x),\tilde{\sigma}(y))=d(x,y)=\|x-y\|.
\]
Therefore,
\[
\|x-y\|=d(\sigma(x),\tilde{\sigma}(y))
\le
\|\sigma(x)-\tilde{\sigma}(y)\|.
\]
Since \(r\mapsto w(B(x,r))\) is increasing, it follows that
\[
w(B(x,\|\sigma(x)-\tilde{\sigma}(y)\|))
\ge
w(B(x,\|x-y\|)).
\]
Hence
\begin{align*}
\|G_{\sigma,\tilde{\sigma}}^{(2)}\|_{L^p(\mu_{\sigma,\tilde{\sigma}})}^p
&=
\int_{\Gamma}\int_{\Gamma}
\frac{|f(\sigma(x))-f(\sigma(y))|^p}
{w(B(x,\|\sigma(x)-\tilde{\sigma}(y)\|))^2}
\,dw(x)\,dw(y)
\\
&\le
\int_{\Gamma}\int_{\Gamma}
\frac{|f(\sigma(x))-f(\sigma(y))|^p}
{w(B(x,\|x-y\|))^2}
\,dw(x)\,dw(y)
\\
&=
\|F_\sigma\|_{L^p(\mu)}^p.
\end{align*}
Thus
\[
\|G_{\sigma,\tilde{\sigma}}^{(2)}\|_{L^p(\mu_{\sigma,\tilde{\sigma}})}
\le
\|F_\sigma\|_{L^p(\mu)}.
\]
Consequently,
\[
\|G_{\sigma,\tilde{\sigma}}^{(1)}\|_{L^p(\mu_{\sigma,\tilde{\sigma}})}
\le
\|F_\sigma\|_{L^p(\mu)}
+
\|G_{\sigma,\tilde{\sigma}}^{(3)}\|_{L^p(\mu_{\sigma,\tilde{\sigma}})}.
\]
Summing over all \(\sigma\neq \tilde{\sigma}\), we obtain
\begin{align*}
&\sum_{\sigma\in G}\|F_\sigma\|_{L^p(\mu)}
+
\sum_{\substack{\sigma,\tilde{\sigma}\in G\\ \sigma\neq \tilde{\sigma}}}
\|G_{\sigma,\tilde{\sigma}}^{(1)}\|_{L^p(\mu_{\sigma,\tilde{\sigma}})}
\\
&\lesssim
\sum_{\sigma\in G}\|F_\sigma\|_{L^p(\mu)}
+
\sum_{\substack{\sigma,\tilde{\sigma}\in G\\ \sigma\neq \tilde{\sigma}}}
\|G_{\sigma,\tilde{\sigma}}^{(3)}\|_{L^p(\mu_{\sigma,\tilde{\sigma}})}.
\end{align*}

For the reverse inequality, note that
\[
G_{\sigma,\tilde{\sigma}}^{(3)}
=
G_{\sigma,\tilde{\sigma}}^{(1)}
-
G_{\sigma,\tilde{\sigma}}^{(2)}.
\]
Hence
\[
\|G_{\sigma,\tilde{\sigma}}^{(3)}\|_{L^p(\mu_{\sigma,\tilde{\sigma}})}
\le
\|G_{\sigma,\tilde{\sigma}}^{(1)}\|_{L^p(\mu_{\sigma,\tilde{\sigma}})}
+
\|G_{\sigma,\tilde{\sigma}}^{(2)}\|_{L^p(\mu_{\sigma,\tilde{\sigma}})}.
\]
Using again
\[
\|G_{\sigma,\tilde{\sigma}}^{(2)}\|_{L^p(\mu_{\sigma,\tilde{\sigma}})}
\le
\|F_\sigma\|_{L^p(\mu)},
\]
we get
\begin{align*}
&\sum_{\sigma\in G}\|F_\sigma\|_{L^p(\mu)}
+
\sum_{\substack{\sigma,\tilde{\sigma}\in G\\ \sigma\neq \tilde{\sigma}}}
\|G_{\sigma,\tilde{\sigma}}^{(3)}\|_{L^p(\mu_{\sigma,\tilde{\sigma}})}
\\
&\le
\sum_{\sigma\in G}\|F_\sigma\|_{L^p(\mu)}
+
\sum_{\substack{\sigma,\tilde{\sigma}\in G\\ \sigma\neq \tilde{\sigma}}}
\|G_{\sigma,\tilde{\sigma}}^{(1)}\|_{L^p(\mu_{\sigma,\tilde{\sigma}})}
+
\sum_{\substack{\sigma,\tilde{\sigma}\in G\\ \sigma\neq \tilde{\sigma}}}
\|F_\sigma\|_{L^p(\mu)}
\\
&\le
|G|
\sum_{\sigma\in G}\|F_\sigma\|_{L^p(\mu)}
+
\sum_{\substack{\sigma,\tilde{\sigma}\in G\\ \sigma\neq \tilde{\sigma}}}
\|G_{\sigma,\tilde{\sigma}}^{(1)}\|_{L^p(\mu_{\sigma,\tilde{\sigma}})}
\\
&\le
|G|
\left(
\sum_{\sigma\in G}\|F_\sigma\|_{L^p(\mu)}
+
\sum_{\substack{\sigma,\tilde{\sigma}\in G\\ \sigma\neq \tilde{\sigma}}}
\|G_{\sigma,\tilde{\sigma}}^{(1)}\|_{L^p(\mu_{\sigma,\tilde{\sigma}})}
\right).
\end{align*}
This proves \eqref{eq:claim-non-dunkl}. Substituting \eqref{eq:claim-non-dunkl} into \eqref{eq:non-dunkl-decomp-2}, we conclude that
\[
\|f\|_{B^p(\mathbb R^N,dw)}
\sim
\sum_{\sigma\in G}\|F_\sigma\|_{L^p(\mu)}
+
\sum_{\substack{\sigma,\tilde{\sigma}\in G\\ \sigma\neq \tilde{\sigma}}}
\|G_{\sigma,\tilde{\sigma}}^{(3)}\|_{L^p(\mu_{\sigma,\tilde{\sigma}})}.
\]
This completes the proof.
\end{proof}

{\color{red} The difficulty for establishing the equivalence of two Besov spaces: In the case of $N=1$ and $R=\{\pm\sqrt{2}\}$, the following Lemma was established in your previous note:
\begin{lemma}
We have
\begin{align*}
\sum_{\substack{\sigma,\tilde{\sigma}\in G\\ \sigma\neq \tilde{\sigma}}}\|F_{\sigma,\tilde{\sigma}}^{(3)}\|_{L^p(\mu)}\sim\sum_{\substack{\sigma,\tilde{\sigma}\in G\\ \sigma\neq \tilde{\sigma}}}
\|G_{\sigma,\tilde{\sigma}}^{(3)}\|_{L^p(\mu_{\sigma,\tilde{\sigma}})}
\end{align*}
\end{lemma}%
%\begin{proof}
%Fix $\sigma,\tilde{\sigma}\in G$ with $\sigma\neq \tilde{\sigma}$. We have
%$$\|F_{\sigma,\tilde{\sigma}}^{(3)}\|_{L^p(\mu)}=\left(\int_{\Gamma}\int_{\Gamma}\left(\frac{\|x-y\|}{\|\sigma(x)-\tilde{\sigma}(y)\|}\right)^{p}\frac{|f(\sigma(y))-f(\tilde{\sigma}(y))|^p}{w(B(x,\|x-y\|))^2}dw(x)dw(y)\right)^{1/p}$$
%Recall that
%$${w(B(x,r))}\sim r^N \prod_{v \in R}(|\langle x, v \rangle |+r)^{k(v)}$$
%$$
%dw(x) :=\prod_{v \in R} |\langle x,v\rangle|^{k(v)} dx,
%$$
%We have
%\begin{align*}
%&\|F_{\sigma,\tilde{\sigma}}^{(3)}\|_{L^p(\mu)}\\
%&\sim\left(\int_{\Gamma}\int_{\Gamma}\left(\frac{\|x-y\|}{\|\sigma(x)-\tilde{\sigma}(y)\|}\right)^{p}\frac{|f(\sigma(y))-f(\tilde{\sigma}(y))|^p}{\|x-y\|^{2N} \prod_{v \in R}(|\langle x, v \rangle |+\|x-y\|)^{2k(v)}}\prod_{v_1 \in R} |\langle x,v_1\rangle|^{k(v_1)}\prod_{v_2 \in R} |\langle y,v_2\rangle|^{k(v_2)}dxdy\right)^{1/p}
%\end{align*}
%
%On the other hand,
%\begin{align*}
%&\|G_{\sigma,\tilde{\sigma}}^{(3)}\|_{L^p(\mu_{\sigma,\tilde{\sigma}})}\\
%&= \int_{\Gamma}\int_{\Gamma}
%\frac{|f(\sigma(y))-f(\tilde{\sigma}(y))|^p}
%{w(B(x,\|\sigma(x)-\tilde{\sigma}(y)\|))^2}
%\,dw(x)\,dw(y)\\
%&\sim \int_{\Gamma}\int_{\Gamma}
%\frac{|f(\sigma(y))-f(\tilde{\sigma}(y))|^p}
%{\|\sigma(x)-\tilde{\sigma}(y)\|^{2N}\prod_{v\in R}(|\langle x, v \rangle |+\|\sigma(x)-\tilde{\sigma}(y)\|)^{2k(v)}}
%\,\prod_{v_1 \in R} |\langle x,v_1\rangle|^{k(v_1)}\prod_{v_2 \in R} |\langle y,v_2\rangle|^{k(v_2)}dxdy\\
%\end{align*}
%
%
%
%\end{proof}
However, the above Lemma fails for general $N$ and $R$, see the following counterexample in Lemma \ref{lem:counterexample-2d-cross-terms}.
}
\begin{lemma}\label{lem:counterexample-2d-cross-terms}
Let
\[
R=\{(\pm \sqrt2,0),(0,\pm \sqrt2)\}\subset \mathbb R^2,
\qquad
G=\{I,s_1,s_2,s_{12}\}\cong \{\pm 1\}^2,
\]
where
\[
s_1(x_1,x_2)=(-x_1,x_2),\qquad
s_2(x_1,x_2)=(x_1,-x_2),\qquad
s_{12}=s_1s_2.
\]
Let \(\Gamma=\mathbb R_+^2\). After renormalizing the measure by a harmless positive constant, we may write
\[
dw(x)=|x_1|^a|x_2|^b\,dx,
\qquad a,b\ge 0.
\]
Assume that \(p>2\) and \(b>p-2\). Then there exists a function
\(f\in L^p_{\mathrm{loc}}(\mathbb R^2,dw)\) such that
\[
\sum_{\substack{\sigma,\tilde{\sigma}\in G\\ \sigma\neq \tilde{\sigma}}}
\|G_{\sigma,\tilde{\sigma}}^{(3)}\|_{L^p(\mu_{\sigma,\tilde{\sigma}})}
<\infty,
\]
but
\[
\sum_{\substack{\sigma,\tilde{\sigma}\in G\\ \sigma\neq \tilde{\sigma}}}
\|F_{\sigma,\tilde{\sigma}}^{(3)}\|_{L^p(\mu)}
=\infty.
\]
Consequently,
\[
\sum_{\substack{\sigma,\tilde{\sigma}\in G\\ \sigma\neq \tilde{\sigma}}}
\|G_{\sigma,\tilde{\sigma}}^{(3)}\|_{L^p(\mu_{\sigma,\tilde{\sigma}})}
\not\sim
\sum_{\substack{\sigma,\tilde{\sigma}\in G\\ \sigma\neq \tilde{\sigma}}}
\|F_{\sigma,\tilde{\sigma}}^{(3)}\|_{L^p(\mu)}.
\]
\end{lemma}

\begin{proof}
Recall that in this setting
\[
d\mu(x,y)=\frac{1}{w(B(x,\|x-y\|))^2}\,dw(x)\,dw(y),
\]
and, for \(\sigma,\tilde{\sigma}\in G\),
\[
d\mu_{\sigma,\tilde{\sigma}}(x,y)
=
\frac{1}{w(B(x,\|\sigma(x)-\tilde{\sigma}(y)\|))^2}\,dw(x)\,dw(y),
\qquad (x,y)\in \Gamma\times\Gamma.
\]
We also recall that
\[
F_{\sigma,\tilde{\sigma}}^{(3)}(x,y)
=
\frac{\|x-y\|}{\|\sigma(x)-\tilde{\sigma}(y)\|}
\bigl(f(\sigma(y))-f(\tilde{\sigma}(y))\bigr),
\]
whereas
\[
G_{\sigma,\tilde{\sigma}}^{(3)}(x,y)
=
f(\sigma(y))-f(\tilde{\sigma}(y)).
\]

We shall construct a function \(f\) for which the two sums behave differently.

Choose a nonzero function \(\eta\in C_c^\infty((1,2))\), and define
\[
\psi(t):=t^{-\gamma}\mathbf 1_{(0,1)}(t),\qquad t>0,
\]
where \(\gamma>0\) is chosen so that
\begin{equation}\label{eq:gamma-choice}
\frac{p-1}{p}<\gamma<\frac{b+1}{p}.
\end{equation}
Such a choice is possible precisely because \(b>p-2\), which is equivalent to
\[
b+1>p-1.
\]

Now define
\[
f(x_1,x_2):=\operatorname{sgn}(x_1)\,\eta(|x_1|)\,\psi(|x_2|),
\qquad (x_1,x_2)\in \mathbb R^2.
\]

We first note that \(f\in L^p_{\mathrm{loc}}(\mathbb R^2,dw)\). Indeed,
\[
|f(x_1,x_2)|^p
\lesssim
\mathbf 1_{\{1<|x_1|<2,\ 0<|x_2|<1\}}\,|x_2|^{-\gamma p},
\]
hence
\[
\int_K |f(x)|^p\,dw(x)
\lesssim
\int_{1<|x_1|<2}\int_0^1 x_2^{\,b-\gamma p}\,dx_2\,dx_1
<\infty
\]
for every compact set \(K\), because \(\gamma p<b+1\) by \eqref{eq:gamma-choice}.

We divide the proof into two steps.

\medskip

\noindent
\textbf{Step 1.}
\emph{The sum involving \(G_{\sigma,\tilde{\sigma}}^{(3)}\) is finite.}

Let
\[
E:=[1,2]\times(0,1)\subset \Gamma.
\]
For \(y=(y_1,y_2)\in E\), we have \(\eta(y_1)\neq 0\) on the support of interest and \(\psi(y_2)=y_2^{-\gamma}\).

Observe that \(f\) is odd in the first variable and even in the second variable. Therefore:

\begin{itemize}
\item if \(\sigma\) and \(\tilde{\sigma}\) have the same sign in the first coordinate, then
\[
f(\sigma(y))-f(\tilde{\sigma}(y))=0
\qquad\text{for all } y\in \Gamma;
\]
\item if \(\sigma\) and \(\tilde{\sigma}\) have opposite signs in the first coordinate, then
\[
|f(\sigma(y))-f(\tilde{\sigma}(y))|
=
2\,\eta(y_1)\psi(y_2),
\qquad y\in E.
\]
\end{itemize}

Hence only those ordered pairs \((\sigma,\tilde{\sigma})\) with opposite first signs can contribute nontrivially, and there are only finitely many such pairs. It therefore suffices to show that, for every such pair,
\[
\|G_{\sigma,\tilde{\sigma}}^{(3)}\|_{L^p(\mu_{\sigma,\tilde{\sigma}})}<\infty.
\]

Fix such a pair \((\sigma,\tilde{\sigma})\). Then there exists \(\varepsilon\in\{\pm 1\}\) such that
\[
\|\sigma(x)-\tilde{\sigma}(y)\|^2
=
(x_1+y_1)^2+(x_2-\varepsilon y_2)^2,
\qquad x,y\in \Gamma.
\]
In particular, for \(y\in E\),
\[
\|\sigma(x)-\tilde{\sigma}(y)\|^2
\ge
(x_1+1)^2+(x_2-\varepsilon y_2)^2.
\]
Using the standard ball estimate
\[
w(B(x,r))\sim r^2(x_1+r)^a(x_2+r)^b,
\qquad x\in \Gamma,\ r>0,
\]
we obtain
\[
\frac{1}{w(B(x,r))^2}\,dw(x)
=
\frac{x_1^a x_2^b}{r^4(x_1+r)^{2a}(x_2+r)^{2b}}\,dx
\lesssim
r^{-4-a-b}\,dx,
\]
because
\[
\frac{x_1^a}{(x_1+r)^{2a}}\le r^{-a},
\qquad
\frac{x_2^b}{(x_2+r)^{2b}}\le r^{-b}.
\]
Therefore, uniformly in \(y\in E\),
\begin{align*}
\int_\Gamma \frac{1}{w(B(x,\|\sigma(x)-\tilde{\sigma}(y)\|))^2}\,dw(x)
&\lesssim
\int_{x_1>0,\ x_2>0}
\frac{dx}
{\bigl((x_1+1)^2+(x_2-\varepsilon y_2)^2\bigr)^{(4+a+b)/2}} \\
&\le
\int_{u>1,\ v>-1}
\frac{du\,dv}{(u^2+v^2)^{(4+a+b)/2}}
<\infty.
\end{align*}
Hence
\begin{align*}
\|G_{\sigma,\tilde{\sigma}}^{(3)}\|_{L^p(\mu_{\sigma,\tilde{\sigma}})}^p
&=
\int_\Gamma\int_\Gamma
\frac{|f(\sigma(y))-f(\tilde{\sigma}(y))|^p}
{w(B(x,\|\sigma(x)-\tilde{\sigma}(y)\|))^2}
\,dw(x)\,dw(y) \\
&\lesssim
\int_E |2\eta(y_1)\psi(y_2)|^p\,dw(y) \\
&\lesssim
\int_1^2 |\eta(y_1)|^p y_1^a\,dy_1
\int_0^1 y_2^{\,b-\gamma p}\,dy_2
<\infty,
\end{align*}
again because \(\gamma p<b+1\).

Since only finitely many ordered pairs \((\sigma,\tilde{\sigma})\) are involved, we conclude that
\[
\sum_{\substack{\sigma,\tilde{\sigma}\in G\\ \sigma\neq \tilde{\sigma}}}
\|G_{\sigma,\tilde{\sigma}}^{(3)}\|_{L^p(\mu_{\sigma,\tilde{\sigma}})}
<\infty.
\]

\medskip

\noindent
\textbf{Step 2.}
\emph{The sum involving \(F_{\sigma,\tilde{\sigma}}^{(3)}\) is infinite.}

It is enough to exhibit a single ordered pair \((\sigma,\tilde{\sigma})\) for which
\[
\|F_{\sigma,\tilde{\sigma}}^{(3)}\|_{L^p(\mu)}=\infty.
\]
We take
\[
\sigma=I,\qquad \tilde{\sigma}=s_1.
\]
For \(y=(y_1,y_2)\in E\), define
\[
J(y):=f(y)-f(s_1(y)).
\]
Since \(f\) is odd in the first variable, we have
\[
J(y)=2\eta(y_1)\psi(y_2),\qquad y\in E.
\]
Hence
\begin{align*}
\|F_{I,s_1}^{(3)}\|_{L^p(\mu)}^p
&=
\int_\Gamma\int_\Gamma
\left(\frac{\|x-y\|}{\|(x_1+y_1,x_2-y_2)\|}\right)^p
\frac{|J(y)|^p}{w(B(x,\|x-y\|))^2}
\,dw(x)\,dw(y).
\end{align*}

Fix \(y=(y_1,y_2)\in E\). Consider the square
\[
Q_y:=\Bigl\{x\in \Gamma:\ |x_1-y_1|<\frac{y_2}{4},\ |x_2-y_2|<\frac{y_2}{4}\Bigr\}.
\]
Since \(y_1\in[1,2]\) and \(0<y_2<1\), every \(x\in Q_y\) satisfies
\[
x_1\sim 1,\qquad x_2\sim y_2,\qquad \|x-y\|\lesssim y_2.
\]
Moreover,
\[
\|(x_1+y_1,x_2-y_2)\|\sim 1
\qquad\text{for all }x\in Q_y.
\]
Indeed,
\[
x_1+y_1\ge \left(y_1-\frac{y_2}{4}\right)+y_1\ge 2-\frac14>\frac74,
\]
while \(x_1+y_1\le 2y_1+\frac{y_2}{4}\le \frac{17}{4}\), and also \(|x_2-y_2|<y_2/4<1/4\).

Let \(d:=\|x-y\|\). For \(x\in Q_y\), the ball estimate gives
\[
w(B(x,d))
\sim d^2(x_1+d)^a(x_2+d)^b.
\]
Since \(x_1+d\lesssim 1\) and \(x_2+d\lesssim y_2\) on \(Q_y\), we obtain
\[
w(B(x,d))\lesssim d^2 y_2^b.
\]
Therefore,
\[
\frac{1}{w(B(x,d))^2}\gtrsim d^{-4}y_2^{-2b}.
\]
On the other hand, since \(x_1\sim 1\) and \(x_2\sim y_2\) on \(Q_y\),
\[
dw(x)=x_1^a x_2^b\,dx\gtrsim y_2^b\,dx.
\]
Combining these estimates, we get on \(Q_y\):
\[
\left(\frac{\|x-y\|}{\|(x_1+y_1,x_2-y_2)\|}\right)^p
\frac{1}{w(B(x,\|x-y\|))^2}\,dw(x)
\gtrsim
d^{p-4}y_2^{-b}\,dx.
\]
Hence
\begin{align*}
\int_\Gamma
\left(\frac{\|x-y\|}{\|(x_1+y_1,x_2-y_2)\|}\right)^p
\frac{1}{w(B(x,\|x-y\|))^2}\,dw(x)
&\ge
c\,y_2^{-b}\int_{Q_y} d^{p-4}\,dx.
\end{align*}

Now \(Q_y\) contains a Euclidean ball centered at \(y\) of radius \(c_0y_2\), where \(c_0>0\) is an absolute constant. Since \(p>2\), the function \(z\mapsto |z|^{p-4}\) is locally integrable in dimension \(2\), and thus
\[
\int_{Q_y} d^{p-4}\,dx
\ge
\int_{B(y,c_0y_2)} \|x-y\|^{p-4}\,dx
\sim
y_2^{\,p-2}.
\]
Therefore,
\[
\int_\Gamma
\left(\frac{\|x-y\|}{\|(x_1+y_1,x_2-y_2)\|}\right)^p
\frac{1}{w(B(x,\|x-y\|))^2}\,dw(x)
\gtrsim
y_2^{\,p-2-b},
\qquad y\in E.
\]

Substituting this into the formula for \(\|F_{I,s_1}^{(3)}\|_{L^p(\mu)}^p\), we obtain
\begin{align*}
\|F_{I,s_1}^{(3)}\|_{L^p(\mu)}^p
&\gtrsim
\int_E |J(y)|^p\, y_2^{\,p-2-b}\,dw(y).
\end{align*}
Since \(dw(y)=y_1^a y_2^b\,dy\) on \(\Gamma\), this becomes
\begin{align*}
\|F_{I,s_1}^{(3)}\|_{L^p(\mu)}^p
&\gtrsim
\int_1^2\int_0^1
|2\eta(y_1)\psi(y_2)|^p\, y_2^{\,p-2-b}\, y_1^a y_2^b\,dy_2\,dy_1 \\
&=
2^p\left(\int_1^2 |\eta(y_1)|^p y_1^a\,dy_1\right)
\left(\int_0^1 y_2^{\,p-2-\gamma p}\,dy_2\right).
\end{align*}
By the choice \(\gamma>(p-1)/p\), we have
\[
p-2-\gamma p<-1,
\]
and therefore
\[
\int_0^1 y_2^{\,p-2-\gamma p}\,dy_2=\infty.
\]
Hence
\[
\|F_{I,s_1}^{(3)}\|_{L^p(\mu)}=\infty.
\]
In particular,
\[
\sum_{\substack{\sigma,\tilde{\sigma}\in G\\ \sigma\neq \tilde{\sigma}}}
\|F_{\sigma,\tilde{\sigma}}^{(3)}\|_{L^p(\mu)}
=\infty.
\]

Combining Step 1 and Step 2, we conclude that
\[
\sum_{\substack{\sigma,\tilde{\sigma}\in G\\ \sigma\neq \tilde{\sigma}}}
\|G_{\sigma,\tilde{\sigma}}^{(3)}\|_{L^p(\mu_{\sigma,\tilde{\sigma}})}
<\infty,
\qquad
\sum_{\substack{\sigma,\tilde{\sigma}\in G\\ \sigma\neq \tilde{\sigma}}}
\|F_{\sigma,\tilde{\sigma}}^{(3)}\|_{L^p(\mu)}
=\infty.
\]
This proves the lemma.
\end{proof}

{\color{red}Comment: However, Lemma \ref{lem:counterexample-2d-cross-terms} does not imply that these two Besov spaces are not equivalent.}




\bibliographystyle{plain}
\bibliography{references}

\end{document}
